\tituloResumo
% Máx. 250 palavras: procedimentos + resultados mais relevantes.

\rasc{Neste trabalho estudou-se o perfil Roncz/Marske-7, usado como substituto do
perfil da raiz da asa do Cirrus SR22, com a teoria dos perfis delgados, o método
dos painéis e uma simulação CFD 2D em
OpenFOAM com o modelo de turbulência $k$--$\omega$ SST. No método dos painéis
escolheram-se $N = 260$ painéis, com um erro de 0,5\,\% no $c_l$ a $5^\circ$
face a 800 painéis, e no CFD uma malha com cerca de 64 mil células, com uma
incerteza relativa de 1,3\,\% no $c_l$ entre malhas. Os ângulos de sustentação
nula previstos pela teoria dos perfis delgados, pelo método dos painéis e pelo
CFD são $-0{,}13^\circ$, $-0{,}39^\circ$ e $-0{,}67^\circ$, e os declives na zona
linear são 6,28, 6,78 e 6,26~rad$^{-1}$. O método dos painéis prevê o maior
declive por incluir a espessura, enquanto no CFD esse efeito é quase anulado
pela viscosidade. O CFD é linear até cerca de $10^\circ$ e o perfil entra em
perda entre $15^\circ$ e $20^\circ$, com o escoamento descolado e não
estacionário a $20^\circ$, o que os métodos potenciais não preveem. As linhas de corrente e as distribuições de $C_p$ dos
painéis e do CFD são semelhantes a $0^\circ$ e $5^\circ$, mas muito diferentes a
$20^\circ$. Os métodos potenciais servem para estimativas rápidas na zona
linear, e o CFD é necessário perto e depois da perda e para o arrasto,
com um custo muito maior.}
